% Use the existing owner-derived source mechanism; do not edit frozen originals.
\OLSADeclareDocumentTextCorrection{OLP-0019}{FA0019-RESTRICTION-CONVENTION}
{\cap A^2$.}
{\cap A^2$؛ در این تعریف، هر دو مؤلفهٔ زوج مرتب باید در مجموعهٔ داده‌شده باشند.}
\OLSADeclareDocumentTextCorrection{OLP-0019}{FA0019-APPLICATION-IMAGE}
{\emph{اعمال} $R$ بر $A$ برابر است با}
{\emph{اعمال} $R$ بر $A$، یعنی گرفتن تصویر مجموعه تحت رابطه، برابر است با}
\OLSADeclareDocumentTextCorrection{OLP-0021}{FA0021-FUNCTION-ASSIGNMENT}
{یک \emph{تابع} $f \colon A \to B$ نگاشتی از هر !!{element}
از~$A$ به یک !!{element} از~$B$ است.}
{یک \emph{تابع} $f \colon A \to B$ به هر !!{element}
از~$A$ دقیقاً یک !!{element} از~$B$ نسبت می‌دهد.}
\OLSADeclareDocumentTextCorrection{OLP-0021}{FA0021-MEMBERS-IDIOM}
{!!{element}s از~$A$، ورودی‌ها}
{اعضای~$A$، ورودی‌ها}
\OLSADeclareDocumentTextCorrection{OLP-0021}{FA0021-DEFINITE-VALUE}
{نامیده می‌شوند، و !!{element} از~$B$}
{نامیده می‌شوند، و آن !!{element} از~$B$}
\OLSADeclareDocumentTextCorrection{OLP-0021}{FA0021-CAPTION-ARTICLE}
{\caption{یک تابع نگاشتی از هر !!{element} از یک مجموعه به
    !!a{element} از مجموعه‌ای دیگر است.}
{\caption{یک تابع هر !!{element} از یک مجموعه را به
    یک !!{element} از مجموعه‌ای دیگر می‌فرستد.}
\OLSADeclareDocumentTextCorrection{OLP-0021}{FA0021-RANGE-ALL-VALUES}
{که از مقادیر~$f$ برای برخی آرگومان‌ها تشکیل می‌شود؛}
{که از همهٔ مقادیر~$f$ برای آرگومان‌های دامنه تشکیل می‌شود؛}
\OLSADeclareDocumentTextCorrection{OLP-0021}{FA0021-NOT-FUNCTIONAL-IDIOM}
{تابع‌وار نیست، زیرا}
{تابع تعریف نمی‌کند، زیرا}
\OLSADeclareDocumentTextCorrection{OLP-0021}{FA0021-SQRT-NONNEGATIVE}
{بازگرداندن تنها ریشهٔ دوم مثبت آن را تابع‌وار کنیم:}
{بازگرداندن تنها ریشهٔ دوم نامنفی، یک تابع تعریف کنیم:}
\OLSADeclareDocumentTextCorrection{OLP-0021}{FA0021-SQRT-SOURCE-NOTE}
{\Nat \to \Real$.}
{\Nat \to \Real$.
\emph{یادداشت ویراستاری:} در متن مبدأ «مثبت» آمده است؛ چون دامنه صفر را هم دربر می‌گیرد، در اینجا «نامنفی» لازم است.}
\OLSADeclareDocumentTextCorrection{OLP-0021}{FA0021-INPUT-VARIABLE}
{با داشتن یک عدد طبیعی~$n$، تابع $g$}
{با داشتن یک عدد طبیعی~$x$، تابع $g$}
\OLSADeclareDocumentTextCorrection{OLP-0021}{FA0021-INPUT-SOURCE-NOTE}
{$x+1$، را خروجی خواهد داد.
\end{ex}}
{$x+1$، را خروجی خواهد داد.
\emph{یادداشت ویراستاری:} نماد ورودی در این جمله با نماد به‌کاررفته در تعریف تابع یکسان شد.
\end{ex}}
\OLSADeclareSetup{OLP-0019}{\olsa_source_correction_install_document_text:}
\OLSADeclareSetup{OLP-0021}{\olsa_source_correction_install_document_text:}
